\subsection{Single-opening games, setup, and proof-engine hypotheses}\label{sec:security14:games}
We work throughout over the \emph{odd-characteristic} field
$\F$, with $N=2^m$, $M=2^{m+s}$, $1\le B\le N$
and the full-relation protocol of
Definition~\ref{def:full:protocol}.
Its basic claim is the evaluation of \emph{one}
preprocessing-bound table $f:I_m\to\F$ at \emph{one}
point $z$ chosen after the first commitment.
The classical security parameter is $\lambda$;
statistical distance is denoted $\Delta$, and
computational distinguishing advantage is denoted
$\operatorname{Adv}^{\rm dist}$.

\begin{definition}[Classical commitment setup]\label{def:security14:setup}
Sample a cyclic prime-order group $G=\langle g\rangle$
of order $P>2$ and
$\alpha\overset{\$}{\leftarrow}\mathbb Z_P^\times$.
Publish $h=g^\alpha$, but withhold $\alpha$ in
the \emph{real} setup. The relation-proof engine
has a separate reference string
$\mathsf{crs}_{\rm rel}$, generated by its
specified setup algorithm.
The field-symbol commitment of
Definition~\ref{def:full:symbol-commit}
uses independent uniform blindings for
\emph{every} encoded bit of every named
hidden variable, including the masked words.
The adversary receives $(G,P,g,h,
\mathsf{crs}_{\rm rel})$.
\end{definition}

\begin{definition}[Evaluation-binding game]\label{def:security14:binding-game}
The cheating prover emits a pre-evaluation
commitment registry, the initial word root,
and an accepted preprocessing relation proof.
Only \emph{after} this transcript is fixed
does the honest verifier sample and reveal
$z\in\F^m$.
The prover then announces $v\in\F$
and completes all masked Sumcheck,
fold-oracle, authentication and relation
proof stages of
Definition~\ref{def:full:protocol}.
The game is won if all verifier checks
accept but
\[
 v\neq\widetilde f_{\rm pre}(z),
\]
where $f_{\rm pre}$ is the single source
table extracted by the pre-evaluation
knowledge extractor.
Extraction failure is counted as a
separate bad event, not as an arbitrary
choice of $f_{\rm pre}$ after observing $z$.
\end{definition}

\begin{definition}[Authorized one-opening zero-knowledge game]\label{def:security14:zk-game}
A classical malicious verifier chooses its
point $z$ \emph{after} receiving the
pre-evaluation commitment and relation proof.
It may choose later field challenges and
authenticated query indices causally from
its complete view. The authorized
functionality
$\mathcal F_{\rm eval}^f$ reveals
$\widetilde f(z)$ to the simulator
when $z$ has been selected.
The full view includes
$\mathsf{crs}_{\rm rel}$, the verifier's
random tape, all commitment vectors,
Merkle roots and paths, all relation-proof
transcripts and the public evaluation
claim $v$. The privacy experiment makes
one opening; no repeated evaluation query
to the same source commitment is included.
\end{definition}

\begin{definition}[Online relation-engine extraction contract]\label{def:security14:extract-contract}
For the fixed finite family of circuit
relations in Chapter~13, let $K$ be the
maximum number of accepted proof invocations
on any opening transcript.
A \emph{prefix-preserving extraction
implementation} of $\Pi_{\rm rel}$ supplies
a reduction with an extraction method for
each accepted circuit proof. Except for
an event of probability at most
$\varepsilon_{\rm KE}$ over the \emph{whole}
adaptive execution, these extractions:
\begin{enumerate}[label=\textup{(\roman*)}]
\item give actual openings of every
registered Pedersen bit commitment and
valid witnesses to every accepted circuit
statement;
\item extract one source table and mask
from the preprocessing proof \emph{before}
the real point $z$ is sampled;
\item complete the extractions of all
pre-challenge registered values before
their respective fresh outer challenges;
\item leave the accepted main transcript
and its future honest outer-challenge
distributions unchanged; and
\item return all witnesses with the same
named public commitment vectors and
Merkle roots seen by the verifier.
\end{enumerate}
For a suitable trapdoor-extractable NIZK
argument of knowledge this can be sought
via its independent extraction key;
for a rewindable interactive argument it
requires a separate checkpoint theorem.
\emph{Chapter~12's one-bit affine extractor
alone does not discharge this contract
for arbitrary hash and field circuits.}
\end{definition}

\begin{definition}[Adaptive witness-indistinguishability contract]\label{def:security14:wi-contract}
The engine $\Pi_{\rm rel}$ has
$\varepsilon_{\rm WI}$-\emph{adaptive
multi-proof witness indistinguishability}
if, for every classical PPT verifier and
every sequence of at most $K$ adaptively
selected \emph{true} public statements,
its entire joint transcript when proofs
use one sequence of valid witnesses is
computationally indistinguishable,
with advantage at most
$\varepsilon_{\rm WI}$, from its
transcript when they use another sequence
of valid witnesses for the same public
statements. The verifier's auxiliary
information, previous proof transcripts,
causal challenge choices and any retained
public commitment roots are included.
This is a \emph{joint} contract; no
marginal-message hiding shortcut
is permitted. If each of $K$ adaptive
proof stages contributes at most
$\delta_{\rm WI}$ under a valid
conditional hybrid, one may set
$\varepsilon_{\rm WI}\le K\delta_{\rm WI}$.
\end{definition}

\begin{definition}[Reduction error terms]\label{def:security14:errors}
Write $\varepsilon_{\rm DL}$ for the
success probability of the concrete
reduction that derives
$\log_g h$ from inconsistent extracted
bit openings, and
$\varepsilon_{\rm CRH}$ for the
success probability of its collision
finder against the Merkle hash.
Write $\varepsilon_{\rm KE}$ for
Definition~\ref{def:security14:extract-contract}
and $\varepsilon_{\rm WI}$ for
Definition~\ref{def:security14:wi-contract}.
These are \emph{assumption-dependent
advantages}, not unconditional
numerical constants.
\end{definition}

\begin{lemma}[Completeness carries to the final security game]\label{lem:security14:complete}
If the relation engine has perfect
completeness and the honest prover
uses the field, masks and query
schedule of Chapter~13, the
full-relation verifier accepts every
correct single evaluation with
probability one.
\end{lemma}
\begin{proof}
The game uses exactly the honest
algorithm and verifier predicates
of Definition~\ref{def:full:protocol}.
Theorem~\ref{thm:full:completeness}
supplies a valid hidden witness
for every statement and proves
all algebraic and Merkle checks.
The stipulated proof-engine
completeness turns those valid
witnesses into accepted proofs.
\end{proof}

\begin{remark}[Not yet an instantiated NP argument]
The commitment group and full arithmetic
circuits are concrete. The contracts
of Definitions~\ref{def:security14:extract-contract}
and~\ref{def:security14:wi-contract}
specify the \emph{precise} remaining
cryptographic properties of the
chosen NP relation-proof engine.
A concrete scheme, CRS, proof format
and reduction meeting both contracts
must be supplied before labeling the
whole construction a deployed secure PCS.
\end{remark}
